\section{傅里叶级数}\label{sec:Fourier_Series}

十九世纪初，傅里叶为了求解热传导问题中的偏微分方程，断言：任何周期函数都可以通过无穷多个正弦和余弦函数的线性组合来表示，即傅里叶级数。这一猜想引发了许多数学家的质疑，促进了分析学的发展，直到1829年，狄利克雷提出了傅里叶级数收敛的充分条件和严格证明。今天，傅里叶级数已经成为数学、物理中的重要工具。

\subsection{三角级数与复指数级数}\label{sub:tri_exp_series}
我们首先回顾数学分析中学过的几个计算傅里叶级数的公式。将周期为T的函数f展开为三角函数形式的傅里叶级数：
\begin{align*}
    f(t) & =\frac{a_0}{2}+\sum_{k = 1}^{\infty} a_k \cos(k\omega t)+b_k\sin(k\omega t) \\
         & =\frac{c_0}{2}+\sum_{k = 1}^{\infty} c_k\cos(k\omega t+\varphi _k)
\end{align*}
（其中$\omega =\frac{2\pi }{T}$为\textbf{基波角频率}，$k\omega (k>1,k\in \mathbb{Z} )$
为k次\textbf{谐波角频率}）
\begin{claim}[三角函数形式的傅里叶系数公式]
\[a_k=\frac{2}{T}\int_T f(t)\cos(k\omega t)\,dt\]
\[b_k=\frac{2}{T}\int_T f(t)\sin(k\omega t)\,dt\]
\[c_k=\sqrt{a_k^2+b_k^2}\]
\end{claim}
下面用完备标准正交函数系的观点来得到以上公式。在学习数学分析时，我们已经看到三角函数系$1,\sin(\omega t),\cos(\omega t),\sin(2\omega t),\cos(2\omega t),\dots\left(\omega =\dfrac{2\pi}{T}\right)$是正交的（读者可以自行验证），但不是单位正交的，因为
\[(\sin(k\omega t),\sin(k\omega t))=\int_{T}\sin^2(k\omega t)\,dt=\int_{T}\frac{1-\cos(2k\omega t)}{2}=\frac{T}{2}\]
\[(\cos(k\omega t),\cos(k\omega t))=\int_{T}\cos^2(k\omega t)\,dt=\int_{T}\frac{1+\cos(2k\omega t)}{2}=\frac{T}{2}\]
可以将它们单位化，也可以直接采用完备正交基分解的公式，
\[a_k=\frac{(f(t),\cos(k\omega t))}{(\cos(k\omega t),\cos(k\omega t))}=\frac{\int_{T}f(t)\cos(k\omega t)^*(t)\,dt}{\int_{T}|\cos(k\omega t)|^2\,dt}
    =\frac{2}{T}\int_{T}f(t)\cos(k\omega t)(t)\,dt\]
\[b_k=\frac{(f(t),\sin(k\omega t))}{(\sin(k\omega t),\sin(k\omega t))}=\frac{\int_{T}f(t)\sin(k\omega t)^*(t)\,dt}{\int_{T}|\sin(k\omega t)|^2\,dt}
    =\frac{2}{T}\int_{T}f(t)\sin(k\omega t)(t)\,dt\]

如果将傅里叶级数展开式写为
\begin{align*}
    f(t) = & \frac{a_0}{2}+\sum_{k = 1}^{\infty} a_k \cos(k\omega t) \\
           & +\sum_{k = 1}^{\infty} b_k \sin(k\omega t)
\end{align*}
则前半部分为偶函数，称之为$f(t)$的\textbf{偶分量}$f_e(t)$；后半部分为奇函数，称之为$f(t)$的\textbf{奇分量}$f_o(t)$。函数的偶分量和奇分量都是唯一的，并且\begin{align*}
    f_e(t)=\frac{f(t)+f(-t)}{2} \\
    f_o(t)=\frac{f(t)-f(-t)}{2}
\end{align*}

容易利用对称性得到：
\begin{claim}[偶函数的傅里叶展开]
当$f(t)$为偶函数，或者由$f(t)$做偶延拓时,展开式为
\[f(t)=\frac{a_0}{2}+\sum_{k = 1}^{\infty} a_k\cos(k\omega t)\]
其中
\[a_k=\frac{4}{T}\int_{0}^{T/2} f(t)\cos(k\omega t)\,dt\]
\[b_k=0\]
\end{claim}\begin{claim}[奇函数的傅里叶展开]
当$f(t)$为奇函数，或者由$f(t)$做奇延拓时,展开式为
\[f(t)=\sum_{k = 1}^{\infty} b_k \sin(k\omega t)\]
其中
\[a_k=0\]
\[b_k=\frac{4}{T}\int_{0}^{T/2} f(t)\sin(k\omega t)\,dt\]
\end{claim}

除了奇偶性，还可以从奇次谐波、偶次谐波的角度来理解函数。
\begin{definition}[奇谐函数]
    函数$f(t)$称为\textbf{奇谐函数}，如果后半个周期的函数是前半个周期的负镜像，即\begin{equation*}
    f\left(t+\frac{T}{2}\right)=-f(t)
\end{equation*}
\end{definition}
奇谐函数的傅里叶级数展开只有奇次谐波分量：
\begin{align*}
    a_k =&\frac{2}{T}\int_{T}f(t)\cos(k\omega t)\,dt\\
    &=\frac{2}{T}\int_{-T/2}^{0}f(t)\cos(k\omega t)\,dt+\frac{2}{T}\int_{0}^{T/2}f(t)\cos(k\omega t)\,dt\\
    &=\frac{2}{T}\int_{0}^{T/2}\left(f(t)\cos(k\omega t)+f(t-T/2)\cos(k\omega(t-T/2))\right)\,dt\\
    &=(1-\cos(k\pi))\frac{2}{T}\int_{0}^{T/2}f(t)\cos(k\omega t)\,dt\\
    &=\begin{cases}
        0,&\text{if }k\text{为偶数}\\
        \frac{4}{T}\int_{0}^{T/2}f(t)\cos(k\omega t)\,dt,&\text{if }k\text{为奇数}
    \end{cases}
\end{align*}
同理有\begin{align*}
    b_k=\begin{cases}
        0,&\text{if }k\text{为偶数}\\
        \frac{4}{T}\int_{0}^{T/2}f(t)\sin(k\omega t)\,dt,&\text{if }k\text{为奇数}
    \end{cases}
\end{align*}
有了奇谐函数自然也能够定义偶谐函数：$f\left(t+\dfrac{T}{2}\right)=f(t)$，容易验证它只有偶次谐波分量，但从定义式可以看出这只是周期减半的函数。

由欧拉公式$\mathrm{e}^{\mathrm{i}k\omega t}=\cos(k\omega t)+\mathrm{i}\sin(k\omega t)$，得到
\[\cos(k\omega t)=\frac{\mathrm{e}^{\mathrm{i}k\omega t}+\mathrm{e}^{-\mathrm{i}k\omega t}}{2},\sin(k\omega t)=\frac{\mathrm{e}^{\mathrm{i}k\omega t}-\mathrm{e}^{-\mathrm{i}k\omega t}}{2\mathrm{i}}\]
故函数$f(t)$也可在指数函数系下展开：\begin{definition}[指数形式傅里叶展开]
    \[f(t)=\sum_{k = 0}^{\infty}  c_k \mathrm{e}^{\mathrm{i}k\omega t} ,c_k=\frac{a_k-\mathrm{i}b_k}{2},c_{-k}=\frac{a_k+\mathrm{i}b_k}{2}=c_k^*,k\in \mathbb{N}\]
（特别地，$c_0=c_0^*\implies c_0\in\mathbb{R} $）
\end{definition}
容易验证$\{\mathrm{e}^{\mathrm{i}k\omega t}\}_{k=0}^{\infty}$是完备正交函数系：
\begin{align*}
    (\mathrm{e}^{\mathrm{i}k_1\omega t},\mathrm{e}^{\mathrm{i}k_2\omega t})&=\int_{T}\mathrm{e}^{\mathrm{i}k_1\omega t}(\mathrm{e}^{\mathrm{i}k_2\omega t})^*\,dt\\
    &=\int_{T}\mathrm{e}^{\mathrm{i}(k_1-k_2)\omega t}\\
    &=\frac{2}{\mathrm{i}\omega (k_1-k_2)}\evalat{\mathrm{e}^{\mathrm{i}(k_1-k_2)\omega t}}{0}{T}\\
    &=0(k_1,k_2\in \mathbb{Z},k_1\neq k_2)
\end{align*}
\begin{align*}
    (\mathrm{e}^{\mathrm{i}k\omega t},\mathrm{e}^{\mathrm{i}k\omega t})&=\int_{T}\mathrm{e}^{\mathrm{i}k\omega t}(\mathrm{e}^{\mathrm{i}k\omega t})^*\,dt\\
    &=\int_{T}\,dt=T(k\in \mathbb{Z} )
\end{align*}
和三角函数系的情况一样，我们得到$c_k$的计算公式：
\begin{align*}
    c_k & =\frac{(f(t),\mathrm{e}^{\mathrm{i}k\omega t})}{(\mathrm{e}^{\mathrm{i}k\omega t},\mathrm{e}^{\mathrm{i}k\omega t})} \\
        & =\frac{1}{T}\int_{T}f(t)(\mathrm{e}^{\mathrm{i}k\omega t})^*\,dt                 \\
        & =\frac{1}{T}\int_{T}f(t)\mathrm{e}^{-\mathrm{i}k\omega t}\,dt
\end{align*}
有时也将$c_k$记作$\hat{f}(k\omega)$或$\hat{F}(k)$，表示f在频域中的点$k\omega$处的值。一般而言，我们只将最小正周期称为一个函数的周期，\textbf{但周期为T的函数可以有多个频率$k\omega(k \in \mathbb{Z})$}。

\subsection{频谱}\label{sub:spectrum}

为了直观地看出周期信号的各个频率分量的大小与相位，可以绘制\textbf{频谱} (spectrum)来表征函数的频域信息。频谱指函数频域图像，通常用其指数形式的傅里叶系数$c_k$来表征。由于难以画出复数，将频谱拆分为\textbf{幅度谱}$|\hat{f}(k\omega)|-\omega$和\textbf{相位谱}$\phi_k-\omega$，其中$\phi_k=\arg \hat{f}(k\omega)$。

对于实信号，
\[\hat{f}(k\omega)=\hat{f}(-k\omega)^*\implies |\hat{f}(k\omega)|=|\hat{f}(-k\omega)|,\phi_{-k}=-\phi_k\]
即幅度谱为偶函数，相位谱为奇函数，所以实信号的频谱中有一半是冗余的。因此，有时也按照三角形式的傅里叶展开式
\[\frac{c_0}{2}+\sum_{k = 1}^{\infty} c_k\cos(k\omega t+\varphi _k)\]
绘制频谱$|c_k|-\omega$（注意此时$c_k$不是指数函数形式的傅里叶系数）和$\phi_k-\omega$，称之为\textbf{单边频谱}，而完整的频谱称为\textbf{双边频谱}，根据欧拉公式，
\begin{align*}
    c_k\cos(k\omega t+\varphi _k)&=\frac{c_k}{2}\mathrm{e}^{\mathrm{i}(k\omega t+\varphi_k)}+\frac{c_k}{2}\mathrm{e}^{-\mathrm{i}(k\omega t+\varphi_k)}\\
    &=\frac{c_k}{2}\mathrm{e}^{\mathrm{i}(k\omega t+\varphi_k)}+\frac{c_{-k}^*}{2}\mathrm{e}^{\mathrm{i}(-k\omega t+\varphi_{-k})}\\
    \implies c_k&=2\hat{f}(k\omega),k>0
\end{align*}
因此\textbf{对于实函数，单边频谱相比双边频谱，在正频率处的幅值加倍，相位不变，展示频域信息的效果不变，但没有冗余}。

\begin{example}[周期矩形脉冲信号的傅里叶级数展开和频谱图]
脉冲宽度为$\tau$，脉冲幅度为E，周期为$T (\tau<T)$的周期矩形脉冲信号$f(t)=E\chi_{[-\tau/2,\tau/2]}$，基波角频率$\omega=\dfrac{2\pi}{T}$，傅里叶级数展开为
\[f(t)=\sum_{k = 1}^{\infty}  \frac{E\tau}{T}\text{Sa}\left(\frac{k\omega \tau}{2}\right)\mathrm{e}^{\mathrm{i}k\omega t}\]
\end{example}
\begin{figure}[H]
    \centering
    \includegraphics[width=0.6\textwidth]{Figure_2}
    \caption{周期矩形脉冲信号及其频谱}
\end{figure}
如图，可以看到，这个频谱与采样函数$\text{Sa}(\omega)$非常相似（为了体现这一点，绘制频谱时将基波角频率大幅减小，并不是第一张图直接做傅里叶级数展开的结果），原因将在\ref{sub:fourier_of_signal}中给出。

\begin{definition}[带宽]
    \textbf{带宽} (bandwidth)指信号的最高频率与最低频率之差，表征信号频率的集中程度。对于实信号，有时仅考虑正频率\cite{signal_systems}，带宽减半。

    周期矩形脉冲信号的频谱是无限的，但能量基本集中在最靠近y轴的两个零点之间，此时可以将它的带宽定义为\textbf{第一过零点带宽}：
    $$B=\dfrac{2\pi}{\tau}$$
\end{definition}

\subsection{傅里叶级数的收敛性}\label{sub:fourier_series_convergency}

我们根据内积空间的理论写出了函数的傅里叶级数展开式，但这并不能保证傅里叶级数存在且逐点收敛于原函数，换言之，前文所述的展开式中使用等号是不严谨的做法，我们只能保证后者是前者的展开式，尽管在多数情况下等号是正确的。

本小节从两个角度来考察傅里叶级数的收敛性：我们将证明傅里叶级数依范数收敛（即均方收敛），并指出一些与逐点收敛性有关的经典结论。

回顾依范数收敛的定义：
\begin{definition}
    在内积空间$L^2([0,T])$中，函数项级数$\sum_{i = 1}^{n}  a_i f_i(t)$依范数收敛于$f(t)$，如果
    \[\lim_{n \to \infty} \left\| f(t)-\sum_{i = 1}^{n}  a_i f_i(t)\right\|_2=0\]
    其中
    \[\left\|f\right\|_2^2=(f,f)=\int_{0}^{T}|f(t)|^2\,dt,\forall f\in L^2([0,T])\]
\end{definition}

作为一个引理，我们首先来证明（指数形式的）傅里叶级数是完备正交基$\{\mathrm{e}^{\mathrm{i}k\omega t}\}_{k=0}^{\infty}$意义下的最佳逼近：
\begin{lemma}
    设$f\in L^2([0,T])$，指数形式的傅里叶系数为$c_k(k\in\mathbb{Z} )$，则对于任意的正整数N和任意的复数列$\{\alpha_k\}_{k=-\infty}^{\infty}$，都有
    \[\left\|f(t)-\sum_{k=-N}^{N}c_k \mathrm{e}^{\mathrm{i}k\omega t}\right\|_2\leq\left\|f(t)-\sum_{k=-N}^{N}\alpha_k \mathrm{e}^{\mathrm{i}k\omega t}\right\|_2\]
    取等条件为$\alpha_k=c_k,k=-N,\cdots,N$。
\end{lemma}
\begin{proof}
    利用正交基的定义，$c_k$是函数$f$在基函数$\mathrm{e}^{\mathrm{i}k\omega t}$上投影分量的大小，因此，
    \[\left(f-\sum_{k=-N}^{N}c_k \mathrm{e}^{\mathrm{i}k\omega t},\mathrm{e}^{\mathrm{i}k\omega t}\right)=0,k=-N,\cdots,N\]
    对下式
    \[f(t)-\sum_{k=-N}^{N}\alpha_k \mathrm{e}^{\mathrm{i}k\omega t}=\left(f(t)-\sum_{k=-N}^{N}c_k \mathrm{e}^{\mathrm{i}k\omega t}\right)+\sum_{k=-N}^{N}(c_k-\alpha_k) \mathrm{e}^{\mathrm{i}k\omega t}\]
    用勾股定理，得到
    \begin{align*}
        \left\|f(t)-\sum_{k=-N}^{N}\alpha_k \mathrm{e}^{\mathrm{i}k\omega t}\right\|_2^2&=\left\|f(t)-\sum_{k=-N}^{N}c_k \mathrm{e}^{\mathrm{i}k\omega t}\right\|_2^2+\left\|\sum_{k=-N}^{N}(c_k-\alpha_k) \mathrm{e}^{\mathrm{i}k\omega t}\right\|_2^2\\
        &\geq \left\|f(t)-\sum_{k=-N}^{N}c_k \mathrm{e}^{\mathrm{i}k\omega t}\right\|_2^2
    \end{align*}
    取等条件为$\alpha_k=c_k,k=-N,\cdots,N$。
\end{proof}

为了证明依范数收敛性，必须引入一个结论：
\begin{claim}
    三角多项式（即形如傅里叶级数部分和的函数）在$L^2([0,T])$空间中是稠密的，即$\forall f\in L^2([0,T]),\forall \varepsilon >0$，存在一个三角多项式$P(t)$使得$\|f-P\|_2<\varepsilon$。
\end{claim}

由于傅里叶级数是$f$的最佳逼近，只要指出一个$M$次的三角多项式，使得$\|f-P\|_2<\varepsilon$，就有
\begin{align*}
    \|f-S_N\|_2 & \leq \|f-P\|_2<\epsilon,\forall N>M
\end{align*}
综上我们证明了：
\begin{theorem}[傅里叶级数的依范数收敛性]
    设$f\in L^2([0,T])$，指数形式的傅里叶系数为$c_k(k\in\mathbb{Z} )$，则
    \[\lim_{N \to \infty} \left\| f(t)-\sum_{k=-N}^{N}c_k \mathrm{e}^{\mathrm{i}k\omega t}\right\|_2=0\]
\end{theorem}

对于逐点收敛性，狄利克雷给出了它的一个充分条件，称之为\textbf{狄利克雷条件}，满足这个条件的函数的傅里叶级数在任意点收敛于其左右极限的平均值。在此之前给出两个概念：

\begin{definition}[狄利克雷条件]
    \quad
    \begin{itemize}[nosep, left=0pt]
    \item $\int_{T}|f(t)|\,dt<\infty$
    \item 在一个周期内f连续或有有限个第一类间断点，即\textbf{分段连续} (piecewise continuous)
    \item 在一个周期内，f的极值点个数有限
\end{itemize}
\end{definition}
\begin{definition}[分段连续]
    函数$f$在区间$[a,b]$上是\textbf{分段连续}的，如果存在区间的划分$a=x_0<x_1<\dots <x_n=b$，使得在每个子区间$(x_{i-1},x_i)$上f连续，且区间端点处$f$的左右极限存在，记作$f\in PC[a,b]$。
\end{definition}

\begin{theorem}[狄利克雷定理]
    满足狄利克雷条件时，f的傅里叶级数展开在在任意点收敛到其左右极限的平均值：
    \[\forall t\in \mathbb{R},\sum_{k=\infty}^{\infty}c_k \mathrm{e}^{\mathrm{i}2\pi kt}=\frac{f(t+)+f(t-)}{2}\]
\end{theorem}

如果函数在有限个点处不连续，我们的确无法将这一结论加强到逐点收敛，因为在不连续点处函数的取值不影响积分值，也就不影响傅里叶系数，当然也就得不到不连续点处的收敛性，但是如果将分段连续函数在不连续点处的取值定义为左右极限的平均值，这样就能够得到逐点收敛性了。

最后一个条件实际上相当于要求f是有界变差函数，感兴趣的读者可以参考实变函数的教材（如\cite{simple_real_analysis}，\cite{real_analysis}）。在\ref{sub:uniform_converge}中我们将讨论另外的更易理解的条件：
\begin{definition}[分段光滑]
    在区间$[a,b]$上，函数$f$称为\textbf{分段光滑} (piecewise smooth)，如果存在区间的划分$a=x_0<x_1<\dots <x_n=b$，使得在每个子区间$(x_{i-1},x_i)$上f可微，且区间端点处$f$的左右导数存在，记作$f\in PS[a,b]$。
\end{definition}
在分段光滑的条件下，我们可以证明傅里叶级数是\textbf{一致收敛}于原函数的。

此外，逐点收敛性可以推广为\textbf{Cesàro收敛}或\textbf{Abel收敛}，逐点收敛性强于这两种收敛性，但是，只要函数$f\in L^1([0,T])\cap L^2([0,T])$，就能保证$f$的傅里叶级数在$f$的连续点处Cesàro收敛于$f(t)$，而Cesàro收敛又能推出Abel收敛。这一条件比狄利克雷条件弱一些。对于这两种收敛性的定义及其证明，可以参考\cite{stein2011fourier}。

\subsection{帕塞瓦尔恒等式}\label{sub:paseval}
类比有限维线性空间中的勾股定理，有
\begin{theorem}[帕塞瓦尔定理/瑞利恒等式]
    \[\|f\|_2=\sum_{k=1}^{\infty}c_k^2|\mathrm{e}^{\mathrm{i}k\omega t}|^2=T\sum_{k=1}^{\infty}c_k^2\]
即\[P=\frac{1}{T}\int_{T}|f(t)|^2\,dt=\sum_{k=1}^{\infty}c_k^2\]
其中P为平均功率。
\end{theorem}
至此，我们提出了傅里叶系数的公式和帕塞瓦尔定理，但所用到的勾股定理是基于有限维线性空间的，现在就来填补这个逻辑漏洞，对此不感兴趣的读者可以忽略这部分内容。设$\{\phi_n\}_{n=1}^{\infty}$是内积空间$L^2(a,b)$的标准正交基，$f\in L^2(a,b)$（注意这里已经不局限于讨论傅里叶级数）。

\begin{lemma}[*贝塞尔不等式]
\[\sum_{n=1}^{\infty}|(f,\phi_n)|^2\leq\|f\|_2^2\]
\end{lemma}
\begin{proof}
    由勾股定理，
    \begin{flalign*}
     \left\|\sum_{n=1}^{N}(f,\phi_n)\phi_n\right\|_2^2= \sum_{n=1}^{N}\left(f,(f,\phi_n)\phi_n\right)=\sum_{n=1}^{N}\overline{(f,\phi_n)}(f,\phi_n)=\sum_{n=1}^{N}|(f,\phi_n)|^2
\end{flalign*}
因此，对任意正整数N，有
\begin{flalign*}
     0&\leq\left\|f-\sum_{n=1}^{N}(f,\phi_n)\phi_n\right\|_2^2\\
     &=\|f\|_2^2-2\text{Re}\ \left(f,\sum_{n=1}^{N}(f,\phi_n)\phi_n\right)+\left\|\sum_{n=1}^{N}(f,\phi_n)\phi_n\right\|_2^2                        \\
     &=\|f\|_2^2-2\sum_{n=1}^{N}|(f,\phi_n)|^2+\sum_{n=1}^{N}|(f,\phi_n)|^2=\|f\|_2^2-\sum_{n=1}^{N}|(f,\phi_n)|^2
\end{flalign*}
令$N\to\infty$即证。从第一行到第二行用到了恒等式$\|\mathbf{a+b}\|^2=\|\mathbf{a}\|^2+2\text{Re}\ \langle\mathbf{a,b}\rangle+\|\mathbf{b}\|^2$，Re表示取实部，这个结论十分简单，留予读者自证。
\end{proof}
在最终的结论帕塞瓦尔定理中这个不等号将变成等号，但它是不可或缺的，并且我们还将在\ref{sub:dirichlet_kernel}中见到它。

\begin{lemma}*
级数$\sum_{n=1}^{N}(f,\phi_n)\phi_n$依范数收敛，并且$\left\|\sum_{n=1}^{\infty}(f,\phi_n)\phi_n\right\|_2\leq\|f\|_2$
\end{lemma}
\begin{proof}
    由贝塞尔不等式，
\begin{flalign*}
     \sum_{n=1}^{\infty}|(f,\phi_n)|^2\leq\|f\|_2^2<\infty,n\to\infty\text{时}|(f,\phi_n)|\to 0
\end{flalign*}
任取$m_1,m_2\in\mathbb{N},m_1<m_2$，由勾股定理，
\begin{flalign*}
              \left\|\sum_{n=m_1}^{m_2}(f,\phi_n)\phi_n\right\|_2^2=\sum_{n=m_1}^{m_2}\left|(f,\phi_n)\right|^2\to 0
\end{flalign*}
因此$\sum_{n=1}^{\infty}(f,\phi_n)\phi_n$构成柯西列。由于$L^2(a,b)$是无限维的完备度量空间，即\textbf{希尔伯特空间} (Hilbert space)，见\ref{sub:Lp_space}。

令$m_1=1,m_2\to\infty$，即得
\begin{flalign*}
\left\|\sum_{n=1}^{\infty}(f,\phi_n)\phi_n\right\|_2 =\sum_{n=1}^{\infty}\left|(f,\phi_n)\phi_n\right|^2\leq\|f\|_2^2
\end{flalign*}
\end{proof}

构建这个引理是为了使用希尔伯特空间中内积的连续性，其表述见下一个命题。

\begin{proposition}[*希尔伯特空间中的内积的连续性]
对于希尔伯特空间$H$，如果$\phi_n\in H$，级数$\sum_{n=1}^{\infty}\phi_n$的部分和$S_N$依范数收敛到S，则任给$y\in H$，总有
\[\lim_{N\to\infty}\left\langle S_n,y\right\rangle=\langle S,y\rangle\]
\end{proposition}
\begin{proof}
\begin{align*}
     & \langle S,y\rangle-\lim_{N\to\infty}\left\langle S_n,y\right\rangle=\lim_{N\to\infty}\langle S-S_n,y\rangle                                     \\
     & \lim_{N\to\infty}\|S-S_N\|=0\implies \lim_{N\to\infty}|\left\langle S-S_n,y\right\rangle|\leq\lim_{N\to\infty}\left\|S-S_N\right\|\|y\|=0               \\
     & \hspace{3cm}\implies \lim_{N\to\infty}\langle S-S_n,y\rangle=0\implies\lim_{N\to\infty}\left\langle S_n,y\right\rangle=\langle S,y\rangle
\end{align*}
\end{proof}

\begin{theorem}*
以下三个命题是等价的：（对于符号$\sim$，参考\ref{sub:Lp_space}）
\begin{enumerate}
    \item $\forall n,(f,\phi_n)=0\implies f\sim 0$，即$\{\phi_n\}_{n=1}^{\infty}$是完备的标准正交基
    \item $\forall f\in L^2(a,b)$，有$f\sim\sum_{n=1}^{\infty}(f,\phi_n)\phi_n$
    \item $\forall f\in L^2(a,b)$，有\textbf{帕塞瓦尔恒等式}：
          \[\|f\|_2^2=\sum_{n=1}^{\infty}|(f,\phi_n)|^2\]
\end{enumerate}
\end{theorem}
\begin{proof}
我们将证明$1\implies 2\implies 3\implies 1$.\\
$1\implies 2$：\begin{align*}
     & \text{令}g\sim \left(f-\sum_{n=1}^{\infty}(f,\phi_n)\phi_n.\right)                                                              \\
     & \forall m\in\mathbb{N},\left(g,\phi_m\right)=\left(f,\phi_m\right)-\sum_{n=1}^{\infty}\left(f,\phi_n\right)\left(\phi_n,\phi_m\right)=\left(f,\phi_m\right)-\left(f,\phi_m\right)=0 \\
\end{align*}
根据1知g=0，即2.这里内积与求和的换序是由命题2.3保证的。\\
$2\implies 3$：由勾股定理，
\[\|f\|_2^2=\lim_{N\to\infty}\left\|\sum_{n=1}^{N}(f,\phi_n)\phi_n\right\|_2^2=\lim_{N\to\infty}\sum_{n=1}^{N}|(f,\phi_n)|^2=\sum_{n=1}^{\infty}|(f,\phi_n)|^2\]
$3\implies 1$：$(f,\phi_n)=0\implies\|f\|_2=0\implies f\sim 0$.
\end{proof}
